% TOP_PROBLEM: {"problemId": "problem.polya-s-conjecture-for-dirichlet-eigenvalues", "canonicalTitle": "Pólya's Conjecture for Dirichlet Eigenvalues", "releaseRank": 402, "primaryDomain": "analysis_pde", "primaryDomainLabel": "Analysis & PDE", "displayStatus": "open_with_solved_subcases", "releaseStatus": "open_with_solved_subcases"}
% !TeX program = lualatex
% Definition and sources only; this file is not an LLM solution attempt.
\documentclass[11pt]{article}
\usepackage[margin=25mm]{geometry}
\usepackage{fontspec}
\setmainfont{Latin Modern Roman}
\usepackage{amsmath,amssymb,mathtools}
\usepackage{unicode-math}
\setmathfont{Latin Modern Math}
\usepackage{xurl,hyperref}
\hypersetup{colorlinks=true,urlcolor=blue}
\setcounter{secnumdepth}{0}
\setlength{\parindent}{0pt}
\setlength{\parskip}{5pt}
\setlength{\emergencystretch}{3em}
\begin{document}
\section*{\texorpdfstring{Pólya's Conjecture for Dirichlet Eigenvalues}{Pólya's Conjecture for Dirichlet Eigenvalues}}\label{402}
\subsection{Definitions and mathematical statement}
Let $\Omega\subset{\mathbb{R}}^d$ be nonempty, bounded and open, $d\ge2$. The Dirichlet Laplacian is the nonnegative self-adjoint operator associated with the quadratic form $u\mapsto\int_\Omega|\nabla u|^2$ on $H_0^1(\Omega)$. List its eigenvalues increasingly with multiplicity as $\lambda_1\le\lambda_2\le\cdots$, and put $\omega_d=|\{x:\|x\|_2<1\}|$. P\'olya's assertion is
\[
 \lambda_k(\Omega)\ge(2\pi)^2
 \left(\frac{k}{\omega_d|\Omega|}\right)^{2/d}
 \qquad(k\ge1).
\]
No smoothness, connectedness, or tiling assumption on $\Omega$ is included. The inequality is for every eigenvalue, not just the leading asymptotic as $k\to\infty$. Neumann eigenvalues require a different formulation and inequality direction.
\par\smallskip\textit{Formulation and concept references:} \href{https://link.springer.com/article/10.1007/s00222-023-01198-1}{[S1]}.

\subsection{Short English statement}
Does every Dirichlet eigenvalue of every bounded Euclidean domain lie above its corresponding Weyl-law value?
\subsection{Sources}
\begin{itemize}
\item[S1]N. Filonov, M. Levitin, I. Polterovich, and D. Sher, Inventiones Mathematicae 234 (2023), 129–169.
\url{https://link.springer.com/article/10.1007/s00222-023-01198-1}.
\textit{Formulation source.}
\end{itemize}
\end{document}
